\chapter{Toward a Hermitian theory of \texorpdfstring{$\mathcal I$}{I}-goodness}
\label{chap:hermitian-igood}

The source book relates three descriptions of a singularity: approximation by analytic singularities, equality of a mass and a multiplier-ideal volume, and equality of two envelopes. The Hermitian problem is to determine when these descriptions remain equivalent. We construct the ideal envelope and compare it with the ordinary envelope without a mass hypothesis. The envelopes of dyadic analytic regularizations then give a precise monotone convergence criterion for their equality. This criterion will be distinguished from a topology on arbitrary singularity types and from a volume characterization.

\section{The source characterization}

Let $X$ be a connected compact K\"ahler manifold of dimension $n$, let $\theta$ be a smooth closed big real $(1,1)$-form, and let $u\in\PSH(X,\theta)$ have positive non-pluripolar mass. Write $\mathrm{MA}_\theta$ for the non-pluripolar Monge--Amp\`ere measure. Use the ordinary envelope $P_\theta[u]$ of \cref{def:singularity-envelope} and the multiplier-ideal envelope $P_\theta[u]_{\mathcal I}$ defined in \cref{def:hermitian-i-envelope}. Let $d_S$ be the pseudometric on singularities from \cite{SGPT}, constructed from the K\"ahler geometry of potentials. The source characterization theorem states the equivalence of
\[
 \begin{gathered}
 u\text{ lies in the }d_S\text{-closure of analytic singularities},\\
 \int_X\mathrm{MA}_\theta(u)
  =\int_X\mathrm{MA}_\theta(P_\theta[u]_{\mathcal I}),\\
 P_\theta[u]=P_\theta[u]_{\mathcal I}.
 \end{gathered}
\]
We recall this as an external input from \cite[Theorem~7.1.1]{SGPT}. We have not constructed such a metric for a general non-closed background.

The middle expression uses the source book's \emph{multiplier-ideal volume}. It is not either of the upper and lower volumes in \cref{def:singularity-volumes}. Those take a supremum or infimum over boundedly equivalent potentials, whereas the source expression first replaces $u$ by an envelope defined using multiplier ideals. We preserve this distinction in the notation below.

\section{An envelope defined by multiplier ideals}

Let $X$ be a connected compact complex manifold of dimension $n\geq1$, and let $\theta$ be a smooth real $(1,1)$-form with $\PSH(X,\theta)\ne\varnothing$. Use the ordinary envelope $P_\theta[u]$ of \cref{def:singularity-envelope} and put $V_\theta=P_\theta(0)$. For quasi-plurisubharmonic functions $u,v$, write $u\preceq v$ if $u\leq v+C$ for some real $C$, and $u\simeq v$ if both inequalities hold. The sheaf $\mathcal I(\lambda u)$ is the multiplier ideal for the weight $\mathrm{e}^{-\lambda u}$, as in \cref{chap:framework}. Write $u\sim_{\mathcal I}v$ when
\[
 \mathcal I(\lambda u)=\mathcal I(\lambda v)
 \quad\text{for every }\lambda>0,
\]
using the weight $\mathrm{e}^{-\lambda u}$ fixed in \cref{chap:framework}.

\begin{definition}\label{def:hermitian-i-envelope}
For $u\in\PSH(X,\theta)$ define
\[
 P_\theta[u]_{\mathcal I}
 :=\left(\sup\{v\in\PSH(X,\theta):v\leq0,
                   \ v\sim_{\mathcal I}u\}\right)^*.
\]
\end{definition}

The family is nonempty because $u-\sup_Xu$ belongs to it; bounded perturbations do not change multiplier ideals. It is bounded above by zero, so its regularized envelope is a genuine $\theta$-plurisubharmonic function. It depends only on the multiplier-ideal equivalence class of $u$. These statements require neither closedness nor a mass bound.

The local ideal calculus of \cref{thm:local-i-inputs,lem:hermitian-i-calculus} shows that taking this supremum preserves every multiplier ideal.

\begin{theorem}\label{thm:hermitian-i-projection}
For every smooth real background $\theta$ and $u\in\PSH(X,\theta)$,
\[
 P_\theta[u]_{\mathcal I}\sim_{\mathcal I}u,\qquad
 P_\theta[P_\theta[u]_{\mathcal I}]_{\mathcal I}=P_\theta[u]_{\mathcal I}.
\]
The regularization in \cref{def:hermitian-i-envelope} does not change the pointwise supremum. Moreover,
\[
 \sup_X P_\theta[u]_{\mathcal I}=0.
\]
\end{theorem}
\begin{proof}
Write $q=P_\theta[u]_{\mathcal I}$. By Choquet's lemma choose a sequence from the defining family whose regularized supremum is $q$. Replace it by successive finite maxima. By \cref{lem:hermitian-i-calculus}, the resulting sequence $v_j$ remains in the family, is increasing, and converges to $q$ almost everywhere.

Fix $x\in X$ and $\lambda>0$. On a coordinate ball about $x$, choose one smooth function $r$ such that every $v_j+r$ and $q+r$ is plurisubharmonic. This is possible because they all have the same smooth curvature lower bound $-\theta$. The increasing-limit assertion of \cref{thm:local-i-inputs} gives
\[
 \mathcal I(\lambda(q+r))_x
   =\mathcal I(\lambda(v_j+r))_x
   =\mathcal I(\lambda u)_x
\]
for large $j$; smooth additions give the second equality. Hence $q\sim_{\mathcal I}u$.

Since $q\leq0$, it now belongs to its own defining family. The pointwise supremum is at least $q$ and at most its regularization $q$, proving the assertion about regularization. The defining families for $P_\theta[q]_{\mathcal I}$ and $P_\theta[u]_{\mathcal I}$ coincide, which proves idempotence. Finally, if $\sup_Xq<0$, the function $q-\sup_Xq$ would be a larger competitor, a contradiction.
\end{proof}

\begin{proposition}\label{prop:p-below-i}
For every $u\in\PSH(X,\theta)$,
\[
 u-\sup_Xu\leq P_\theta[u]\leq P_\theta[u]_{\mathcal I}\leq V_\theta.
\]
\end{proposition}
\begin{proof}
Only the middle inequality needs explanation. Let $h\leq0$ be $\theta$-plurisubharmonic with $h\preceq u$, and set
\[
 h'=\max(h,u-\sup_Xu).
\]
It is a normalized competitor below zero, dominates $h$, and differs from $u$ by a bounded amount. Hence $h'\sim_{\mathcal I}u$, so $h\leq P_\theta[u]_{\mathcal I}$. Take the regularized supremum over $h$. The last inequality holds because every competitor for $P_\theta[u]_{\mathcal I}$ is a competitor for $V_\theta$.
\end{proof}

\begin{corollary}\label{cor:ordinary-preserves-i}
For $u\in\PSH(X,\theta)$,
\[
 P_\theta[u]\sim_{\mathcal I}u,\qquad
 P_\theta[P_\theta[u]]_{\mathcal I}=P_\theta[u]_{\mathcal I},\qquad
 P_\theta[P_\theta[u]_{\mathcal I}]=P_\theta[u]_{\mathcal I}.
\]
\end{corollary}
\begin{proof}
The inequalities of \cref{prop:p-below-i} and monotonicity of local integrability give, for every $\lambda>0$,
\[
 \mathcal I(\lambda u)
 \subseteq\mathcal I(\lambda P_\theta[u])
 \subseteq\mathcal I(\lambda P_\theta[u]_{\mathcal I})
 =\mathcal I(\lambda u).
\]
This proves the first assertion, and equality of the defining families proves the second. For $q=P_\theta[u]_{\mathcal I}$, the inequalities $q\leq P_\theta[q]\leq P_\theta[q]_{\mathcal I}=q$ prove the last assertion; the first inequality uses $q\leq0$.
\end{proof}

\begin{corollary}\label{cor:i-envelope-finite}
If $\theta$ is semipositive and big, then $P_\theta[u]_{\mathcal I}$ has finite extended Monge--Amp\`ere mass for every $u\in\PSH(X,\theta)$, without bounded mass on $X$.
\end{corollary}
\begin{proof}
The last identity in \cref{cor:ordinary-preserves-i} makes $P_\theta[u]_{\mathcal I}$ an ordinary envelope fixed point. Apply the fixed-point mass estimate in \cref{prop:fixed-point-finite}.
\end{proof}

The finiteness statement concerns the envelope. It does not imply finiteness for every function with the same multiplier ideals, as \cref{prop:goodness-infinite-mass} below demonstrates.

\begin{proposition}\label{prop:i-envelope-operations}
For $u,v\in\PSH(X,\theta)$, write $u\preceq_{\mathcal I}v$ if $\mathcal I(\lambda u)\subseteq\mathcal I(\lambda v)$ for every $\lambda>0$. Then
\[
 u\preceq_{\mathcal I}v\quad\Longrightarrow\quad
 P_\theta[u]_{\mathcal I}\leq P_\theta[v]_{\mathcal I}.
\]
For $a>0$ one has $P_{a\theta}[a u]_{\mathcal I}=aP_\theta[u]_{\mathcal I}$. For smooth real $(1,1)$-forms $\theta_i$ and potentials $u_i\in\PSH(X,\theta_i)$, $i=1,2$,
\[
 P_{\theta_1}[u_1]_{\mathcal I}+P_{\theta_2}[u_2]_{\mathcal I}
 \leq P_{\theta_1+\theta_2}[u_1+u_2]_{\mathcal I}.
\]
\end{proposition}
\begin{proof}
If $h$ is a competitor for $P_\theta[u]_{\mathcal I}$, the reverse Lelong inequalities in \cref{thm:local-i-inputs} show that $\max(h,v-\sup_Xv)$ is multiplier-ideal equivalent to $v$. It is below zero and dominates $h$. Taking the supremum proves order preservation. Scaling gives a bijection of the two defining families. For addition, the sum of the two envelopes is below zero, is $(\theta_1+\theta_2)$-plurisubharmonic, and is multiplier-ideal equivalent to $u_1+u_2$ by projection and \cref{lem:hermitian-i-calculus}. It is therefore an admissible competitor on the right.
\end{proof}

\Needspace{9\baselineskip}
\begin{proposition}\label{prop:i-envelope-gauge}
Let $\theta$ be a smooth real background, $u\in\PSH(X,\theta)$, and $f\in C^\infty(X,\mathbb R)$. Set $\theta'=\theta+\ddc f$, $u'=u-f$, $m=\min_Xf$ and $M=\max_Xf$. Then
\[
 P_\theta[u]_{\mathcal I}+m
 \leq P_{\theta'}[u']_{\mathcal I}+f
 \leq P_\theta[u]_{\mathcal I}+M.
\]
The identical inequalities hold for ordinary envelopes.
\end{proposition}
\begin{proof}
If $v$ is a competitor for $P_\theta[u]_{\mathcal I}$, then $v-f+m$ is $\theta'$-plurisubharmonic, below zero, and multiplier-ideal equivalent to $u'$. Therefore
$v\leq P_{\theta'}[u']_{\mathcal I}+f-m$.
Taking the envelope gives the first inequality. Conversely, a competitor $v'$ in the new background gives $v'+f-M$ in the old one, which proves the second inequality. For $P$, use the ordinary singularity order instead of multiplier-ideal equivalence in the same calculation.
\end{proof}

Thus changing the smooth representative gives envelopes whose associated currents have equivalent ordinary singularities in the same $\ddc$-class. Literal equality of the normalized potentials is not the gauge-invariant statement.

\section{Why the raw multiplier-ideal volume changes}

Let $X$ be a connected compact complex manifold of dimension $n\geq1$, and let $\theta$ be a smooth real $(1,1)$-form with $\PSH(X,\theta)\ne\varnothing$. Use the ordinary envelope $P_\theta[u]$ of \cref{def:singularity-envelope} and put $V_\theta=P_\theta(0)$. For quasi-plurisubharmonic functions $u,v$, write $u\preceq v$ if $u\leq v+C$ for some real $C$, and $u\simeq v$ if both inequalities hold. The sheaf $\mathcal I(\lambda u)$ is the multiplier ideal for the weight $\mathrm{e}^{-\lambda u}$, as in \cref{chap:framework}. Use the ideal envelope $P_\theta[u]_{\mathcal I}$ of \cref{def:hermitian-i-envelope}; write $u\sim_{\mathcal I}v$ if their multiplier ideals agree at every positive weight, and $u\preceq_{\mathcal I}v$ if $\mathcal I(\lambda u)\subseteq\mathcal I(\lambda v)$ for every $\lambda>0$. The notation $\mathrm{MA}_\theta$ denotes the extended non-pluripolar measure of \cref{chap:products}, with infinite total mass allowed. Fix $u\in\PSH(X,\theta)$.

One might try to copy the source definition by putting
\[
 W_\theta(u)=\int_X\mathrm{MA}_\theta(P_\theta[u]_{\mathcal I}),
\]
whenever the integral is defined. In a non-closed background this quantity need not depend only on the current $\theta+\ddc u$.

\begin{proposition}\label{prop:raw-i-volume-gauge-failure}
There are two smooth decompositions of the same positive current for which the corresponding values of $W$ are different.
\end{proposition}
\begin{proof}
Use the Hermitian form $\omega$ and smooth function $g=u_t$ from \cref{ex:torus-mass}, choosing $t\ne0$ sufficiently small. Both $\omega$ and $\omega'=\omega+\ddc g$ are Hermitian, and
\[
 \omega+\ddc0=\omega'+\ddc(-g).
\]
For a smooth potential in a Hermitian background, the multiplier ideals at all weights are trivial and zero is an admissible competitor for its multiplier-ideal envelope. That envelope is therefore zero. Consequently
\[
 W_\omega(0)=\int_X\omega^2,\qquad
 W_{\omega'}(-g)=\int_X(\omega+\ddc g)^2.
\]
The two integrals differ by the explicit formula \eqref{eq:torus-mass}.
\end{proof}

This obstruction persists on a complex manifold that is itself K\"ahler, and involves only smooth functions. The source proof of invariance uses closed-background mass invariance after identifying equivalent singularities; \cref{prop:i-envelope-gauge} supplies only the latter step.

The comparison with finite upper and lower volumes under bounded mass is developed in \cref{chap:mass-bound-volumes}.

\section{Envelope equality outside bounded mass}

Let $X$ be a connected compact complex manifold of dimension $n\geq1$, and let $\theta$ be a smooth real $(1,1)$-form with $\PSH(X,\theta)\ne\varnothing$. Use the ordinary envelope $P_\theta[u]$ of \cref{def:singularity-envelope} and put $V_\theta=P_\theta(0)$. For quasi-plurisubharmonic functions $u,v$, write $u\preceq v$ if $u\leq v+C$ for some real $C$, and $u\simeq v$ if both inequalities hold. The sheaf $\mathcal I(\lambda u)$ is the multiplier ideal for the weight $\mathrm{e}^{-\lambda u}$, as in \cref{chap:framework}. Use the ideal envelope $P_\theta[u]_{\mathcal I}$ of \cref{def:hermitian-i-envelope}; write $u\sim_{\mathcal I}v$ if their multiplier ideals agree at every positive weight, and $u\preceq_{\mathcal I}v$ if $\mathcal I(\lambda u)\subseteq\mathcal I(\lambda v)$ for every $\lambda>0$. The notation $\mathrm{MA}_\theta$ denotes the extended non-pluripolar measure of \cref{chap:products}, with infinite total mass allowed. Fix $u\in\PSH(X,\theta)$. For a Hermitian form $\omega$, the class $\mathcal E_{\mathrm{env}}(X,\omega,0)$ is the relative envelope class of \cref{def:relative-envelope-class}.

For the moment call
\begin{equation}\label{eq:background-goodness}
 P_\theta[u]=P_\theta[u]_{\mathcal I}
\end{equation}
\emph{envelope goodness in the fixed background}. This is a working definition. We have not proved independence under addition of an arbitrary positive background, nor equivalence with a closure of analytic singularities. Those are part of what makes the source notion more than a formal equality.

\begin{proposition}\label{prop:goodness-infinite-mass}
If a Hermitian metric $\omega$ fails bounded mass, there exists $U\in\PSH(X,\omega)$ such that
\[
 P_\omega[U]=P_\omega[U]_{\mathcal I}=0,
 \qquad U\in\mathcal E_{\mathrm{env}}(X,\omega,0),
 \qquad \mathrm{MA}^{\mathrm{ext}}_\omega(U)(X)=+\infty.
\]
\end{proposition}
\begin{proof}
Take $U$ from \cref{prop:infinite-envelope-class}. Its ordinary envelope is zero and its multiplier ideals are trivial at every weight. Zero is consequently a competitor for $P_\omega[U]_{\mathcal I}$, and all such competitors are below zero. Thus $P_\omega[U]_{\mathcal I}=0$. The other properties have already been proved.
\end{proof}

For this potential, the raw envelope mass $W_\omega(U)=\int_X\omega^n$ is finite, while the mass of $U$ is infinite. Therefore the envelope equality in \eqref{eq:background-goodness} cannot imply the source raw mass-volume identity on the whole Hermitian potential space. The example supplies a concrete reason to impose a finite-mass regime or to replace the invariant before formulating an equivalence theorem.

\Needspace{8\baselineskip}
Positive divisorial currents supply another explicit regime for envelope equality. \Cref{cor:divisorial-current-ordinary-type} constructs a canonical ordinary type from a finite-mass divisorial current and a smooth residual form, and proves $P_\Omega[v]=P_\Omega[v]_{\mathcal I}\simeq v$ in every Hermitian background admitting that representative. The current may have infinitely many components. Exact curve-fibration trace types are examples; \cref{prop:exact-trace-infinite-representatives} still gives representatives of the same ideal class with different finiteness of Monge--Amp\`ere mass.

\section{An \texorpdfstring{$L^1$}{L1} distance on ideal types and its completion}
\label{sec:ideal-lone-completion}

The ideal envelope gives a canonical representative of each multiplier-ideal class in a fixed background. It therefore supplies a simple metric. Its completion, however, need not consist of ideal envelopes. We prove this explicitly on a curve. The example separates a topology available from compactness alone from the singularity geometry sought in the source characterization.

Fix a smooth positive probability volume form $\mathrm{d}V$ and a smooth real background $\theta$ with $\PSH(X,\theta)\ne\varnothing$. Let
\[
 \mathfrak S_{\mathcal I}(X,\theta)
       =\PSH(X,\theta)/{\sim_{\mathcal I}},\qquad
 \mathfrak M_{\mathcal I}(X,\theta)
       =\{q\in\PSH(X,\theta):Q_\theta[q]=q\}.
\]
By \cref{thm:hermitian-i-projection}, the assignment $[u]\mapsto Q_\theta[u]$ is a bijection from the first set to the second. Every member of the second set has supremum zero.

\begin{proposition}\label{prop:ideal-lone-metric}
The formula
\[
 d_{\mathcal I,\theta,\mathrm{d}V}([u],[v])
      =\int_X|Q_\theta[u]-Q_\theta[v]|\,\mathrm{d}V
\]
defines a finite metric on $\mathfrak S_{\mathcal I}(X,\theta)$. This metric space is totally bounded. Its metric completion is canonically the $L^1$ closure of $\mathfrak M_{\mathcal I}(X,\theta)$ inside the compact space $\mathcal K_\theta$ of \cref{thm:normalized-compactness}. Changing $\mathrm{d}V$ to another smooth positive probability volume gives a bi-Lipschitz equivalent metric.
\end{proposition}
\begin{proof}
The envelopes are integrable and depend only on the ideal classes. Symmetry and the triangle inequality are those of the $L^1$ norm. If the distance is zero, the two quasi-plurisubharmonic canonical representatives agree everywhere, since they agree almost everywhere. Projection then shows that the ideal classes agree. The converse is immediate.

The normalized compactness theorem places $\mathfrak M_{\mathcal I}$ inside compact $\mathcal K_\theta$. Every subset of a compact metric space is totally bounded, and its completion is its closure in that space. This proves the asserted completion without any closedness claim for $\mathfrak M_{\mathcal I}$ itself. Finally the ratio of two smooth positive volume forms on compact $X$ is bounded above and below by positive constants. Multiplying the integral by these bounds proves the last assertion.
\end{proof}

For a Hermitian background, \cref{cor:fixed-hermitian-analytic-density} will prove that analytic classes are dense in this metric. The curve example below describes the still larger metric completion. The background is fixed in this construction. The proposition does not assert invariance of this topology under a smooth change of background, or identify it with the $d_S$ topology. More decisively, its completion may add many different functions with the same multiplier ideals.

\Needspace{12\baselineskip}
\begin{theorem}\label{thm:curve-ideal-lone-completion}
Let $Y$ be a compact connected Riemann surface and $\theta$ a smooth positive form of total mass one. Then the normalized ideal envelopes with finitely many logarithmic poles and zero non-pluripolar mass are dense in $\mathcal K_\theta$ for $L^1$. The logarithmic coefficients can be chosen rational, so these approximants have analytic singularities in the convention of \cite{SGPT}. Consequently the completion in \cref{prop:ideal-lone-metric} is all of $\mathcal K_\theta$. The metric on ideal classes is not complete.
\end{theorem}
\begin{proof}
We first consider a finitely supported probability measure
\[
 \nu=\sum_{j=1}^N a_j\delta_{p_j},\qquad a_j>0,\quad\sum_j a_j=1.
\]
The elementary Green-kernel construction on the curve, stated in \cref{thm:curve-trace-realization}, gives a potential $q_\nu$ with
\[
 \theta+\ddc q_\nu=\nu,\qquad\sup_Yq_\nu=0.
\]
For completeness, choose a kernel with $\ddc_xG(x,p)=\delta_p-\theta$. The sum $\sum_j a_jG(x,p_j)$ has this curvature, and subtracting its finite supremum gives $q_\nu$. Near $p_j$ it is $a_j\log|z-p_j|^2$ plus a smooth function; it is smooth away from the finite support. All of its curvature lies on its pole set, so its non-pluripolar measure is zero.

We claim $Q_\theta[q_\nu]=q_\nu$. If $w\sim_{\mathcal I}q_\nu$, the local ideal criterion gives $\nu(w,p_j)=a_j$. On a curve, this Lelong number is exactly the atomic mass of the positive measure $\theta+\ddc w$. That measure has total mass one by distributional Stokes, so its prescribed atoms already exhaust the mass. It follows that $\theta+\ddc w=\nu$. Thus $\ddc(w-q_\nu)=0$. The difference is a distributional harmonic function, hence smooth and constant on the compact connected curve. Among these translates the greatest one below zero is $q_\nu$. This proves the claim. It also shows that every such $q_\nu$ is an ordinary model potential by \cref{prop:p-below-i}.

Now take any $u\in\mathcal K_\theta$. Its curvature $\nu=\theta+\ddc u$ is a positive probability measure. Choose finitely supported probability measures $\nu_j$ converging weakly to $\nu$. One elementary construction partitions a compact metric space into finitely many Borel sets of diameter tending to zero, places their masses at chosen points, and omits the sets of mass zero. Uniform continuity of each continuous test function proves weak convergence.

We may make these weights rational. For a partition with weights $a_1,\ldots,a_N$, choose a large integer $M$, replace $a_i$ by $b_i=\lfloor Ma_i\rfloor/M$ for $i<N$, and put $b_N=1-\sum_{i<N}b_i$. Discard zero weights. The total absolute change is at most $2(N-1)/M$, which can be made to tend to zero along the partitions. This preserves weak convergence. Writing $b_i=m_i/M$ gives the analytic type
\[
 \left(\frac1M,\mathcal O_Y\left(-\sum_i m_i p_i\right)\right)
\]
with smooth local remainders. Arbitrary real weights were sufficient for the preceding construction; rational weights prove the additional assertion about analytic singularities.

Let $q_j=q_{\nu_j}$. By normalized compactness, every subsequence has an $L^1$-convergent further subsequence with limit $q\in\mathcal K_\theta$. Passing to distributions in the curvature identity gives
\[
 \theta+\ddc q=\nu=\theta+\ddc u.
\]
The difference is constant by the same harmonic argument. Since both suprema are zero, it is zero. Every cluster limit is therefore $u$, so the full sequence converges to $u$ in $L^1$. This proves the density assertion and identifies the completion.

To see incompleteness concretely, choose a nonconstant smooth real function $f$. For sufficiently small $\varepsilon>0$, $\theta+\varepsilon\ddc f>0$. Put $u=\varepsilon f-\sup_Y(\varepsilon f)$. It is smooth, normalized and nonconstant. Its multiplier ideals are all trivial, so $Q_\theta[u]=0$, because zero is an admissible competitor. In particular $u\notin\mathfrak M_{\mathcal I}(Y,\theta)$. The preceding approximants $q_j$ are ideal envelopes, and their classes are Cauchy for $d_{\mathcal I,\theta,\mathrm{d}V}$ because $q_j\to u$ in $L^1$. If they converged in that metric to an ideal class, its envelope would agree almost everywhere, and hence everywhere, with $u$. This contradicts $Q_\theta[u]\ne u$.
\end{proof}

In the last example, every approximating non-pluripolar measure is zero, while the smooth limit has the positive measure $\theta+\ddc u$. There is no failure of weak convergence for the positive currents: their atomic measures converge precisely to that smooth measure. The discontinuity occurs when the pole sets are removed. The sequence lies at zero non-pluripolar mass, so it does not contradict the positive-mass hypotheses in the source book's characterization theorems.

This example also proves that $Q_\theta$ is not continuous for $L^1$ on the normalized potentials: $Q_\theta[q_j]=q_j\to u$, but $Q_\theta[u]=0$. Compactifying the ideal classes by this elementary metric adds all normalized potentials on the curve. Those new points cannot be identified merely by their multiplier ideals, since many distinct smooth potentials have the same trivial ideals. A Hermitian singularity theory must specify whether it seeks such a compactification, a topology restricted by positive mass, or a different completion retaining additional data.


\section{Monotone limits of analytic envelopes}
\label{sec:analytic-envelope-limits}

Demailly regularization preserves a curvature lower bound up to a vanishing error. Along dyadic indices, subadditivity also orders its analytic singularity types. The corresponding envelopes therefore decrease. Divisorial estimates identify the limit with the ideal envelope, by the same ideal-theoretic mechanism used in \cite[Theorem~7.1.1, proof of $(3)\Rightarrow(1)$, Step~1]{SGPT}.

Fix $u\in\PSH(X,\theta)$ and a Hermitian form $h$ on $X$. Choose the regularization from \cref{thm:demailly-hermitian-input} along $m_j=2^j$, starting at a sufficiently large index. Denote it by $a_j$ and choose numbers $\epsilon_j>0$ decreasing to zero such that
\begin{equation}\label{eq:dyadic-analytic-backgrounds}
 a_j\in\PSH(X,\theta+\epsilon_jh),\qquad
 a_j\geq u,\qquad
 a_j\text{ has type }(1/m_j,\mathcal I(m_ju)).
\end{equation}
The metric $h$ remains fixed, and the chosen regularization satisfies $a_j\to u$ in $L^1$. The values $a_j$ need not decrease. Their ordinary singularity types do: subadditivity in \eqref{eq:trace-local-ideal-inputs} gives
\[
 \mathcal I(2m_ju)\subseteq\mathcal I(m_ju)^2,
 \qquad a_{j+1}\preceq a_j.
\]
Taking logarithms of generators proves the ordinary comparison, with a constant allowed to depend on $j$.

\begin{theorem}\label{thm:dyadic-analytic-envelope-limit}
For every $u\in\PSH(X,\theta)$ and every choice in \eqref{eq:dyadic-analytic-backgrounds}, put
\[
 q_j=P_{\theta+\epsilon_jh}[a_j].
\]
Then $q_j$ has analytic singularities, and
\[
 q_j=P_{\theta+\epsilon_jh}[a_j]_{\mathcal I}
       \searrow P_\theta[u]_{\mathcal I}
\]
pointwise and in $L^1$. The limit is independent of the regularizing representatives, the admissible error sequence and the metric $h$.
\end{theorem}
\begin{proof}
By \cref{prop:analytic-envelope-type}, $q_j=P_{\theta+\epsilon_jh}[a_j]_{\mathcal I}\sim a_j$. Since the backgrounds decrease and $a_{j+1}\preceq a_j$, the function $q_{j+1}$ is a competitor for $P_{\theta+\epsilon_jh}[a_j]$. Hence $q_{j+1}\leq q_j$.

Put $p=P_\theta[u]_{\mathcal I}$. The inequality $a_j\geq u$ implies $u\preceq_{\mathcal I}a_j$. Regard $p$ in the larger background and use ideal projection and order preservation there. This gives
\[
 u-\sup_Xu\leq p\leq
 P_{\theta+\epsilon_jh}[p]_{\mathcal I}\leq q_j\leq0.
\]
Thus the decreasing limit $q$ is genuine, has supremum zero and belongs to $\PSH(X,\theta)$: the current inequalities pass to the limit, and a nested intersection of the compact maximum sets preserves the supremum. The displayed lower bound also proves $L^1$ convergence by dominated convergence.

\Needspace{6\baselineskip}
Fix a local divisor $E$. The inequalities $u-\sup_Xu\leq q\leq q_j$ imply
\[
 \nu_E(u)\geq\nu_E(q)\geq\nu_E(q_j)
       =\frac1{m_j}\operatorname{ord}_E\mathcal I(m_ju).
\]
By \cref{lem:bergman-divisorial-error}, the last term tends to $\nu_E(u)$. Therefore $q$ and $u$ have the same Lelong number on every local divisor over $X$. This gives the pointwise data on every local modification as well: blow up the point and use that its Lelong number is the generic number on the exceptional divisor; on a curve the point already is a divisor. The local ideal criterion \cref{thm:local-i-inputs} now gives $q\sim_{\mathcal I}u$.

Since $q\leq0$ is $\theta$-plurisubharmonic, it is a competitor for $p=P_\theta[u]_{\mathcal I}$. The earlier inequality $p\leq q$ proves equality. This identifies the limit solely from $\theta$ and the ideal class of $u$, proving every asserted independence.
\end{proof}

\Needspace{17\baselineskip}
\begin{corollary}\label{cor:dyadic-goodness-criterion}
For the sequence above, the following are equivalent:
\[
 P_\theta[u]=P_\theta[u]_{\mathcal I},\qquad
 q_j\longrightarrow P_\theta[u]\text{ in }L^1,
 \qquad q_j\searrow P_\theta[u]\text{ pointwise}.
\]
They hold for one admissible choice of approximation if and only if they hold for every such choice. For the fixed volume $\mathrm{d}V$,
\[
 \lim_j\int_X(q_j-P_\theta[u])\,\mathrm{d}V
   =\int_X(P_\theta[u]_{\mathcal I}-P_\theta[u])\,\mathrm{d}V\geq0,
\]
and equality to zero is another equivalent condition.
\end{corollary}
\begin{proof}
The theorem identifies the pointwise and $L^1$ limit. Equality almost everywhere is equality everywhere by \cref{lem:qpsh-almost-everywhere-order}. Finally $0\leq q_j-P_\theta[u]\leq-(u-\sup_Xu)$ almost everywhere, so dominated convergence gives the integral formula. A nonnegative difference has zero integral precisely when the two canonical representatives agree.
\end{proof}

\subsection{Absorbing the curvature error}

For a Hermitian background $\omega$, take $\theta=h=\omega$ in \eqref{eq:dyadic-analytic-backgrounds}. A rational rescaling absorbs the curvature error and preserves the rational coefficients of analytic singularities. It gives an approximation entirely within $\PSH(X,\omega)$.

\Needspace{14\baselineskip}
\begin{corollary}\label{cor:fixed-hermitian-analytic-density}
Let $\omega$ be Hermitian and let $u\in\PSH(X,\omega)$. There are analytic $b_j\in\PSH(X,\omega)$ with $b_j\to u$ in $L^1$, with decreasing ordinary singularity types, such that
\[
 P_\omega[b_j]=P_\omega[b_j]_{\mathcal I}\searrow P_\omega[u]_{\mathcal I}.
\]
Consequently analytic singularity classes are dense in the entire ideal-type metric space of \cref{prop:ideal-lone-metric}. The equality $P_\omega[u]=P_\omega[u]_{\mathcal I}$ is equivalent to convergence of these ordinary envelopes to $P_\omega[u]$ in $L^1$.
\end{corollary}
\begin{proof}
Use $h=\omega$ in \eqref{eq:dyadic-analytic-backgrounds}. Enlarge the error bounds to positive rational numbers $\epsilon_j\downarrow0$, and set
\[
 c_j=\frac1{1+\epsilon_j},\qquad b_j=c_ja_j.
\]
Then $c_j\uparrow1$ is rational and
\[
 \omega+\ddc b_j
 =c_j\bigl((1+\epsilon_j)\omega+\ddc a_j\bigr)\geq0.
\]
The analytic coefficient is $c_j/m_j$. Also $b_j\to u$ in $L^1$, since $a_j\to u$ there and the sequence has bounded $L^1$ norm. Because $a_{j+1}\preceq a_j$, $c_{j+1}\geq c_j$, and each $a_j$ is bounded above, one has $b_{j+1}\preceq b_j$.

Put $q_j=P_{(1+\epsilon_j)\omega}[a_j]=P_{(1+\epsilon_j)\omega}[a_j]_{\mathcal I}$. The preceding theorem gives $q_j\searrow P_\omega[u]_{\mathcal I}$ pointwise and in $L^1$. Since $b_j$ is analytic, \cref{prop:analytic-envelope-type} and the homogeneity in \cref{prop:i-envelope-operations} give
\[
 r_j:=P_\omega[b_j]=P_\omega[b_j]_{\mathcal I}=c_jq_j.
\]
Here $q_j\leq0$ decreases and $c_j\uparrow1$, so $r_j$ also decreases, with pointwise limit $P_\omega[u]_{\mathcal I}$. Moreover
\[
 P_\omega[u]_{\mathcal I}\leq q_j\leq r_j\leq0.
\]
The genuine lower bound proves $L^1$ convergence by dominated convergence.

The metric distance from $[b_j]$ to $[u]$ is exactly the $L^1$ distance from $r_j$ to $P_\omega[u]_{\mathcal I}$, so it tends to zero. The last assertion follows from the identified limit.
\end{proof}


\Needspace{18\baselineskip}
\section{Approximation in a Hermitian background}

Let $X$ be a connected compact complex manifold and let $\omega$ be a Hermitian form on $X$. For $u\in\PSH(X,\omega)$, use the ordinary and ideal envelopes from \cref{def:singularity-envelope,def:hermitian-i-envelope}.

\Needspace{10\baselineskip}
\begin{problem}\label{qu:i-good-no-bmp}
Construct a topology on ordinary singularity types in arbitrary Hermitian backgrounds for which the closure of analytic singularities is exactly the set satisfying $P_\omega[u]=P_\omega[u]_{\mathcal I}$. It should recover the dyadic-envelope criterion of \cref{cor:dyadic-goodness-criterion,cor:fixed-hermitian-analytic-density}, remain meaningful for infinite extended mass, and specify how convergence behaves under smooth changes of representative and background.
\end{problem}

The dyadic criterion identifies the limit of a specific ideal-controlled approximation; it does not yet define convergence for arbitrary families of ordinary singularity types. By \cref{cor:fixed-hermitian-analytic-density}, the elementary ideal metric makes analytic classes dense everywhere. The finite damped measures of \cref{chap:damped-measures} provide additional analytic data, but do not by themselves characterize multiplier ideals. Their capacity-continuity theorem also shows a limitation: on a normalized family, convergence in capacity already implies convergence of these damped measures. Merely appending damped weak convergence to capacity convergence adds no condition that could distinguish multiplier-ideal approximation. A convergence theory for arbitrary families must retain additional information about singularities, for example controlled ideal inclusions at each fixed weight. Constructing that theory remains open.
